% Part:sets-functions-relations
% Chapter: size-of-sets
% Section: enumerations-alt

\documentclass[../../../include/open-logic-section]{subfiles}

\begin{document}

\olfileid{sfr}{siz}{enm-alt}

\olsection{Opsommingen en \usetoken{S}{enumerable} verzamelingen}

\begin{editorial}
  In deze paragraaf is een opsomming een bijectie met (een eerste deel
  van) $\Nat$. Zo gebeurt het in de verzamelingenleer. Dit wijkt een
  beetje af van de definities in \olref[enm]{sec} en herhaalt alle
  voorbeelden uit die paragraaf. De tekst is hier ook wat korter.
\end{editorial}

Een eindige verzameling kun je eenvoudig vastleggen door haar
!!{element}s op te schrijven. Dat doen we zo:
\[
  A = \{a_1, a_2, \ldots, a_n\}.
\]
Als de !!{element}s $a_1$, \dots, $a_n$ allemaal verschillend zijn,
geeft deze lijst !!a{bijection} tussen $A$ en de eerste $n$
natuurlijke getallen $0$, \dots, $n-1$. Andersom heeft elke eindige
verzameling maar eindig veel !!{element}s. Je kunt de elementen dus
altijd zo aan getallen koppelen. Is $A$ eindig, dan is er met andere
woorden !!a{bijection} tussen $A$ en $\{0, \dots, n-1\}$. Hier is $n$
het aantal !!{element}s van~$A$.

Als we ook bepaalde soorten oneindige verzamelingen toelaten, kunnen
we sommige oneindige verzamelingen op deze manier in een lijst zetten.
We leggen precies vast wat dat betekent: zo'n oneindige verzameling
wordt opgesomd door !!a{bijection} tussen die verzameling en
heel~$\Nat$.

\begin{defn}[Opsomming in de verzamelingenleer]
Een \emph{opsomming} van een verzameling $A$ is !!a{bijection}. Het
bereik---de werkelijk verkregen waarden---is $A$. Het domein, de
toegestane invoerverzameling, is een eerste stuk van de natuurlijke
getallen $\{0, 1, \ldots, n\}$ {of} de hele verzameling natuurlijke
getallen~$\Nat$.
\end{defn}

\begin{explain}
Waarom heet dit een \emph{opsomming}? Als we een verzameling $A$
hebben opgesomd, is er !!a{bijection} $f$ waarmee we de elementen van
$A$ één voor één kunnen tellen. Het $0$-de element is $f(0)$, het
eerste is $f(1)$, \ldots het $n$-de is $f(n)$\ldots.\footnote{Ja, we
beginnen bij $0$. We zouden natuurlijk ook bij~$1$ kunnen beginnen.
Dat maakt weinig verschil. We hoeven dan alleen~$\Nat$ te vervangen
door~$\PosInt$.} De volgende definitie maakt nog duidelijker waarom:
\end{explain}

\begin{defn}
  \ollabel{defn:enumerable}
  Een verzameling~$A$ is !!{enumerable} precies wanneer $A =
  \emptyset$ of wanneer~$A$ een opsomming heeft. We zeggen dat $A$
  !!{nonenumerable} is precies wanneer $A$ niet !!{enumerable} is.
\end{defn}

\begin{explain}
Een verzameling is dus !!{enumerable} als zij leeg is of als je haar
!!{element}s met een opsomming één voor één kunt tellen.
\end{explain}

\begin{ex}
De identiteitsfunctie $\Id{\Nat} \colon \Nat \to \Nat$, met
$\Id{\Nat}(n) = n$, somt gewoon de natuurlijke getallen op. De functie
die de \emph{positieve} natuurlijke getallen $\Nat^+ = \Nat \setminus
\{0\}$ opsomt, is $g(n) = n + 1$: de opvolgerfunctie.
\end{ex}

\begin{prob}
Laat zien dat een verzameling $A$ !!{enumerable} is precies wanneer
$A = \emptyset$ of er !!a{surjection} $f\colon \Nat \to A$ is. Laat
ook zien dat $A$ !!{enumerable} is precies wanneer er !!a{injection}
$g\colon A \to \Nat$ is.
\end{prob}

\begin{ex}
Neem de functies $f\colon \Nat \to \Nat$ en $g \colon \Nat \to \Nat$
die zijn vastgelegd door
\begin{align*}
  f(n) & = 2n \text{ en}\\
  g(n) & = 2n+1
\end{align*}
De eerste somt de even natuurlijke getallen op en de tweede de oneven
natuurlijke getallen. Geen van beide is !!{surjective}. Daarom is
geen van beide een opsomming van $\Nat$.
\end{ex}

\begin{prob}
Geef een opsomming van de kwadraatgetallen $1$, $4$, $9$, $16$, \dots
\end{prob}

\begin{ex}
Laat $\lceil x \rceil$ de \emph{plafondfunctie} zijn. Deze rondt $x$
naar boven af tot het dichtstbijzijnde gehele getal. De volgende
functie $f \colon \Nat \to \Int$,
\[
  f(n) = (-1)^{n} \left\lceil\tfrac{n}{2}\right\rceil
\]
zet de gehele getallen~$\Int$ in deze volgorde:
\[
\begin{array}{c c c c c c c c}
f(0) & f(1) & f(2) & f(3) & f(4) & f(5) & f(6) & \dots \\ \\
\big\lceil \tfrac{0}{2} \big\rceil & -\big\lceil \tfrac{1}{2}\big\rceil &  \big\lceil \tfrac{2}{2} \big\rceil & -\big\lceil \tfrac{3}{2} \big\rceil & \big\lceil \tfrac{4}{2} \big\rceil  & -\big\lceil \tfrac{5}{2}\big\rceil & \big\lceil \tfrac{6}{2} \big\rceil & \dots \\ \\
0 & -1 & 1 & -2 & 2 & -3 & 3& \dots
\end{array}
\]
Je ziet dat $f$ de waarden van $\Int$ één voor één geeft door steeds
heen en weer te ‘springen’ tussen positieve en negatieve gehele getallen. Je kunt
$f$ ook met twee gevallen vastleggen:
\[
f(n) = \begin{cases}
  \frac{n}{2} & \text{als $n$ even is}\\
  -\frac{n+1}{2} & \text{als $n$ oneven is}
  \end{cases}
\]
\end{ex}

\begin{prob}
Laat zien dat als $A$ en $B$ !!{enumerable} zijn, ook $A \cup B$ dat
is.
\end{prob}

\begin{prob}
Laat met inductie naar $n$ zien dat als $A_1$, $A_2$, \dots, $A_n$
allemaal !!{enumerable} zijn, ook $A_1 \cup \dots \cup A_n$ dat is.
\end{prob}
\end{document}